% ===================================================================
%  BAGAN No. 6 — Bagian Dalam Rumah, Perabot, Makanan, Penerangan.
%
%  Traduction, faite en 2026, en indonesien d'aujourd'hui, sur l'ido
%  de texto/io. Regles de la colonne : voir 10-bagan-01.tex. Cinq
%  pieces, cinq numerotations : les cles portent c1 a c5. Le bloc
%  t06-c1-06-1 est l'un des trois ecarts declares de renvoji.py —
%  l'ido y ouvre sur « (6) », qui est le savon, la ou le francais
%  donne « (16) », la femme de chambre ; la colonne ecrit (16).
%
%  guling (25) EST LE MEILLEUR APPARIEMENT D'OBJET DE TOUT LE RELEVE.
%  L'ido dit « transversa kuseno » et le francais « traversin » : un
%  long coussin cylindrique qu'on met en travers du lit. L'indonesien
%  dit guling, et c'est exactement le meme objet — au point que les
%  Neerlandais des Indes l'ont emporte chez eux sous le nom de « Dutch
%  wife ». Ce n'est ni un emprunt ni une periphrase : deux peuples ont
%  le meme coussin et chacun a son mot. Le releve n'avait pas encore
%  rencontre de coincidence aussi franche sur un objet aussi
%  particulier.
%
%  ET teko (31) CONFIRME LA QUATRIEME COUCHE DES LE PREMIER TABLEAU
%  QUI LA SUIT. Le tableau 5 avait revele le hokkien avec « cat » et
%  « loteng » ; ici la theiere est teko, de 茶壺, et elle est dans une
%  cuisine ou tout le reste est malais ou neerlandais. Trois couches
%  se partagent cette seule piece : wajan la poele et talenan la
%  planche a hacher sont malaises, kompor le fourneau et panci la
%  casserole sont neerlandais, teko est hokkien. Une cuisine
%  indonesienne raconte trois siecles d'histoire par ses ustensiles.
%
%  ONZIEME REFUS, ET C'EST LA SALLE DE BAINS. La piece centrale d'une
%  salle de bains indonesienne est le bak mandi : un bassin macconne
%  ou l'on puise avec un gayung pour se verser l'eau dessus — on ne
%  s'y baigne PAS. Ecrire « bak mandi » pour la baignoire (1) aurait
%  mis une salle de bains de Java sur une planche belge, et c'eut ete
%  le remplacement le plus discret du releve, parce que les deux
%  objets sont dans la meme piece et servent au meme geste. La colonne
%  ecrit « bak rendam », le bassin ou l'on se trempe, qui decrit la
%  chose du fac-simile. En revanche handuk la serviette (4) reste :
%  c'est le meme objet, et le mot vient du neerlandais handdoek.
%
%  DEUX CENT ONZE RENVOIS, CINQUANTE-DEUX PAIRES : le tableau le plus
%  charge du volume, et le plus lourd en paires apres le 5. Ce n'est
%  pas un inventaire pur — il attache sans arret un objet a sa piece,
%  a son meuble, a sa proprietaire — et le compte le dit, comme au
%  tableau 13 du volume.
% ===================================================================

%%K t06-apar-1 apar
\VUsaut{6.0mm}
\VUtitre{15.0pt}{60}{BAGAN No. 6}
\VUsaut{1.4mm}
\VUtitre{11.4pt}{0}{Bagian Dalam Rumah, Perabot, Makanan,}
\VUsaut{0.4mm}
\VUtitre{11.4pt}{0}{Penerangan.}
\VUsaut{2.2mm}

%%K t06-apar-2 apar
Rumah itu sudah selesai. Paman Ioannes sudah menempati rumah barunya,
yang tata ruangnya kita kenal. Sekarang mari kita lihat bagaimana ia
memerabotinya. Pertama-tama kita akan menjelajahi :

%%K t06-tit-1 sub
\VUpk{10.2pt}{40}{Kamar tidur.}
\VUsaut{3.0mm}

%%K t06-c1-01-1 p
1. --- Kita datang pada saat yang kurang tepat, sebab pelayan kamar
sedang sibuk membersihkan \VUgras{kamar} itu.

%%K t06-c1-02-1 p
2. --- Di atas \VUgras{meja rias} \textsuperscript{(1)} tersebar
bermacam-macam benda untuk mencuci muka, menyisir rambut, pendeknya
untuk berdandan. Benda-benda itu adalah \VUgras{baskom}
\textsuperscript{(2)} dan \VUgras{kendi air} \textsuperscript{(3)},
\VUgras{sikat rambut} \textsuperscript{(4)}, \VUgras{sisir}
\textsuperscript{(5)}, \VUgras{sabun} \textsuperscript{(6)} dan
\VUgras{spons} \textsuperscript{(7)}, akhirnya sebuah \VUgras{botol
kecil} \textsuperscript{(8)} untuk obat gigi cair.

%%K t06-c1-03-1 p
3. --- Di lantai, di depan meja rias, inilah \VUgras{ember}
\textsuperscript{(9)} untuk menampung air kotor.

%%K t06-c1-04-1 p
4. --- Di \VUgras{ruang muka} \textsuperscript{(10)} kita melihat
beberapa \VUgras{gantungan baju} \textsuperscript{(13)} dan dua
\VUgras{payung} \textsuperscript{(11)} di dalam \VUgras{tempat payung}
\textsuperscript{(12)}.

%%K t06-c1-05-1 p
5. --- Di sisi lain pintu berdiri \VUgras{lemari cermin}
\textsuperscript{(14)} yang berisi \VUgras{kain-kain}
\textsuperscript{(15)}.

%%K t06-c1-06-1 p
6. --- Mari kita manfaatkan saat \VUgras{pelayan kamar}
\textsuperscript{(16)} merapikan \VUgras{tempat tidur}
\textsuperscript{(17)} untuk menjelajahi bagian-bagiannya.
\VUgras{Rangka tempat tidur} \textsuperscript{(20)} berisi
\VUgras{kasur pegas} \textsuperscript{(21)} yang baik, yang di atasnya
Iustina sebentar lagi akan meletakkan \VUgras{kasur}
\textsuperscript{(22)} wol. Kemudian ia akan mengambil dari kursi kedua
\VUgras{seprai} \textsuperscript{(23)} dan \VUgras{selimut}
\textsuperscript{(24)}. Lalu ia akan meletakkan di kepala tempat tidur
\VUgras{guling} \textsuperscript{(25)} dan \VUgras{bantal kepala}
\textsuperscript{(26)}, yang bersama \VUgras{selimut bulu}
\textsuperscript{(27)} kini ada di atas \VUgras{permadani tempat tidur}
\textsuperscript{(32)}. Ia memakai kemoceng untuk membersihkan debu
\VUgras{tirai} \textsuperscript{(19)} dan \VUgras{kelambu}
\textsuperscript{(18)}, dan inilah sebagian pekerjaannya selesai.

%%K t06-c1-07-1 p
7. --- Setelah itu ia harus sedikit merapikan \VUgras{meja kecil}
\textsuperscript{(28)}, memutar kembali \VUgras{jam weker}
\textsuperscript{(29)}, menyiapkan \VUgras{lampu malam}
\textsuperscript{(30)}, dan mengangkat \VUgras{lilin}
\textsuperscript{(31)} untuk membersihkan kaki lilinnya.

%%K t06-tit-2 sub
\VUsaut{5.0mm}
\VUpk{10.2pt}{40}{Keranjang jahit dan perkakas perapian.}
\VUsaut{3.0mm}

%%K t06-c1-08-1 p
8. --- \VUgras{Keranjang jahit} \textsuperscript{(34)} di dekat
\VUgras{bangku kecil} \textsuperscript{(33)} pun sangat perlu
dirapikan. Ia harus melipat \VUgras{sulaman}
\textsuperscript{(35)}, menyimpan di dalam \VUgras{meja jahit}
\textsuperscript{(38)} \VUgras{bidal,} \VUgras{benang}
\textsuperscript{(36)}, \VUgras{jarum} \textsuperscript{(37)} dan
\VUgras{gunting} \textsuperscript{(39)}, lalu mengunci tutup meja itu
dengan \VUgras{kunci} \textsuperscript{(40)}.

%%K t06-c1-09-1 p
9. --- Di atas \VUgras{meja berkaki satu} \textsuperscript{(41)} di
tengah ruangan, Iustina akan mengganti air \VUgras{botol air}
\textsuperscript{(42)}. Ia akan menaruh \VUgras{gula} ke dalam
\VUgras{tempat gula} \textsuperscript{(43)}, membersihkan
\VUgras{penjepit gula} \textsuperscript{(44)} dan mengangkat
\VUgras{sikat baju} \textsuperscript{(46)}.

%%K t06-c1-10-1 p
10. --- Sisi kanan kamar itu menuntut lebih sedikit pekerjaan
darinya, sebab \VUgras{kursi panjang} \textsuperscript{(47)} dan
\VUgras{bantalnya} \textsuperscript{(48)} sudah pada tempatnya, dan
karena cuaca tidak dingin, tidak ada satu pun perkakas
\VUgras{perapian} \textsuperscript{(49)} yang bergeser.
\VUgras{Sekop} \textsuperscript{(50)}, \VUgras{jepitan}
\textsuperscript{(51)}, \VUgras{penyangga kayu} \textsuperscript{(55)}
dan \VUgras{sekat abu} \textsuperscript{(56)} bersih ; \VUgras{sekat
api} \textsuperscript{(54)} tidak dipindahkan, dan \VUgras{peti kayu}
\textsuperscript{(52)} penuh berisi \VUgras{kayu bakar}
\textsuperscript{(53)}.

%%K t06-noto-1 noto
(*) Babo = anak kecil ; I. baby, P. bébé, It. bambino.

%%K t06-noto-2 noto
(*) Lambrequino = P. lambrequin, Sp. lambrequines, Port. lambrequins,
bahkan J. I. lambrequins ; jadi J. I. P. Port. Sp.

%%K t06-c1-11-1 p
11. --- Namun ia masih harus membersihkan debu \VUgras{jam bandul}
\textsuperscript{(57)}, \VUgras{kaki lilin bercabang}
\textsuperscript{(58)} dan \VUgras{tempat barang saku}
\textsuperscript{(59)}, akhirnya mengelap \VUgras{cermin}
\textsuperscript{(60)}.

%%K t06-c1-12-1 p
12. --- Adapun \VUgras{ayunan bayi} \textsuperscript{(61)} si kecil
(si babo (*)), ia sudah siap.

%%K t06-tit-3 sub
\VUsaut{5.0mm}
\VUpk{10.2pt}{40}{Ruang tamu.}
\VUsaut{3.0mm}

%%K t06-c2-01-1 p
1. --- Sekarang inilah ruang tamu. Ia \VUgras{kamar} yang sangat
indah. Perapian, yang terletak di sudut kanan, dihiasi dengan
\VUgras{patung kecil} \textsuperscript{(1)} yang halus, dua
\VUgras{jambangan} \textsuperscript{(3)} yang indah, dan ditutupi
\VUgras{lambrekin} \textsuperscript{(2)} beledu (*).

%%K t06-c2-02-1 p
2. --- Supaya bunga api dari tungku tidak membakar permadani, di
depan tungku dipasang \VUgras{sekat bunga api} \textsuperscript{(4)}
dari tembaga.

%%K t06-c2-03-1 p
3. --- Di dekatnya, sebuah \VUgras{meja tulis} \textsuperscript{(5)}
perempuan yang anggun terbuka. Di atasnya \VUgras{nyonya rumah}
\textsuperscript{(23)} menulis surat-suratnya.

%%K t06-c2-04-1 p
4. --- Jendela itu dihiasi \VUgras{tirai kaca} \textsuperscript{(6)}
dengan \VUgras{tali pengikat} \textsuperscript{(8)} yang disangkutkan
pada \VUgras{cantelan} \textsuperscript{(7)}, dan tirai kain besar yang
\VUgras{berjuntai di atas} \textsuperscript{(9)}.

%%K t06-c2-05-1 p
5. --- Di ceruk jendela, sebuah \VUgras{tempat pot}
\textsuperscript{(11)} berisi \VUgras{tanaman} \textsuperscript{(10)}
hijau, sebatang palem yang bagus, yang daunnya terentang hampir sampai
ke \VUgras{lampu minyak} \textsuperscript{(12)}.

%%K t06-c2-06-1 p
6. --- \VUgras{Kap lampu} \textsuperscript{(13)} pada lampu itu diatur
oleh Maria, \VUgras{gadis} \textsuperscript{(14)} yang memesona, yang
duduk di piano \textsuperscript{(15)}. Dengan duduk sangat tegak di
atas \VUgras{bangku piano} \textsuperscript{(16)}, ia mempelajari
sebuah lagu. Ia sangat cakap dan cermat, sebagaimana dibuktikan oleh
\VUgras{lemari musiknya} \textsuperscript{(17)} yang tertata sempurna.

%%K t06-tit-4 sub
\VUsaut{5.0mm}
\VUpk{10.2pt}{40}{Si anak nakal.}
\VUsaut{3.0mm}

%%K t06-c2-07-1 p
7. --- Saudara laki-lakinya, Paulus, mencari sebuah \VUgras{buku}
\textsuperscript{(19)} di \VUgras{rak buku} \textsuperscript{(18)}. Ia
seorang \VUgras{anak muda} \textsuperscript{(20)} yang tidak bisa diam
dan tak tertahankan. Sambil bergantung pada rak buku untuk meraih buku
yang terlalu tinggi, ia berisiko menjatuhkan \VUgras{patung dada}
\textsuperscript{(21)} yang ada di atasnya.

%%K t06-c2-08-1 p
8. --- Ia memanfaatkan untuk berbuat kebodohan itu saat ibunya sibuk
berbincang dengan seorang teman, seorang \VUgras{nyonya}
\textsuperscript{(22)} yang datang menginap beberapa hari di rumahnya.
Ibunya duduk di \VUgras{sofa} \textsuperscript{(24)} dekat
\VUgras{kursi besar} \textsuperscript{(25)}, dan temannya duduk di
\VUgras{kursi} \textsuperscript{(27)} yang hanya kelihatan
\VUgras{sandarannya} \textsuperscript{(26)}.

%%K t06-c2-09-1 p
9. --- \VUgras{Bingkai} \textsuperscript{(30)} besar yang tergantung
di dinding berisi \VUgras{potret} \textsuperscript{(29)} seorang
kerabat perempuan. Di dalam \VUgras{album} \textsuperscript{(32)} yang
diletakkan di atas meja terkumpul \VUgras{foto-foto}
\textsuperscript{(33)} anggota keluarga yang lain. Anak-anak baru saja
membolak-baliknya, sebab \VUgras{kursi-kursi} \textsuperscript{(31)}
masih berkumpul di sekeliling meja.

%%K t06-c2-10-1 p
10. --- \VUgras{Jambangan Tiongkok} \textsuperscript{(34)} dari
porselen, yang ada di dekat \VUgras{pisau kertas}
\textsuperscript{(35)}, dihadiahkan kepada Maria pada hari
perayaannya, dan \VUgras{sekat lipat} \textsuperscript{(36)} yang
menghiasi sudut kamar dibawa dari Tiongkok oleh pamannya.

%%K t06-c2-11-1 p
11. --- Diriangkan oleh \VUgras{lukisan-lukisan} \textsuperscript{(37)}
yang indah tergantung di \VUgras{dinding} \textsuperscript{(42)},
diterangi oleh \VUgras{lampu gantung} \textsuperscript{(38)} yang
\VUgras{lampu-lampu listriknya} \textsuperscript{(39)} bersinar riang
pada malam hari, ruang tamu itu tampak sangat menyenangkan.
\VUgras{Permadani} \textsuperscript{(40)} lembut yang terbentang di
atas lantai papan, dan \VUgras{bel listrik} \textsuperscript{(41)} yang
begitu mudah dipakai, melengkapi kenyamanannya.

%%K t06-tit-5 sub
\VUsaut{3.6mm}
\VUpk{10.2pt}{40}{Ruang makan.}
\VUsaut{3.0mm}

%%K t06-c3-01-1 p
1. --- Lonceng makan malam telah berbunyi. Seluruh keluarga, kecuali
Maria, berkumpul di sekeliling meja. Ruang makan itu dihangatkan
dengan menyenangkan oleh \VUgras{tungku} \textsuperscript{(1)} porselen
yang sangat baik, dengan \VUgras{cerobong} \textsuperscript{(2)} yang
langsing.

%%K t06-c3-02-1 p
2. --- Kamar itu tidak berisi banyak perabot. Di sepanjang dinding
tergantung, di atas \VUgras{bangku} \textsuperscript{(4)}, sebuah
\VUgras{pancuran kecil} \textsuperscript{(3)} penghias. Sangat dekat
dari situ ada \VUgras{bufet} \textsuperscript{(5)}, tempat orang
meletakkan makanan dingin, hidangan penutup, dan bermacam-macam alat
meja yang belum segera diperlukan.

%%K t06-c3-03-1 p
3. --- Demikianlah di papan atas kita melihat \VUgras{ayam panggang}
\textsuperscript{(7)}, \VUgras{ikan} \textsuperscript{(8)} dan sebuah
\VUgras{mangkuk} \textsuperscript{(10)} berisi \VUgras{buah}
\textsuperscript{(9)}. Di bawahnya inilah \VUgras{keju}
\textsuperscript{(11)}, \VUgras{roti} \textsuperscript{(12)},
\VUgras{tempat botol bumbu} \textsuperscript{(13)} yang berisi segala
yang diperlukan untuk membumbui salad : minyak, cuka, \VUgras{garam}
\textsuperscript{(14)} dan \VUgras{merica} \textsuperscript{(15)}.
Akhirnya, di papan bawah ada \VUgras{kue} \textsuperscript{(17)} di
dekat \VUgras{tempat mostar} \textsuperscript{(16)}.

%%K t06-c3-04-1 p
4. --- Jam ruang makan, yang disebut \VUgras{jam dinding}
\textsuperscript{(20)}, mempunyai bentuk yang sangat khusus. Ia
terpasang di dalam bingkai kayu yang juga berisi \VUgras{barometer}
\textsuperscript{(19)} dan \VUgras{termometer} \textsuperscript{(18)}.

%%K t06-tit-6 sub
\VUsaut{4.6mm}
\VUpk{10.2pt}{40}{Penyajian makan malam.}
\VUsaut{3.0mm}

%%K t06-c3-05-1 p
5. --- Pada semua dinding kamar itu dipasang \VUgras{panel kayu}
\textsuperscript{(21)} dari kayu ek yang dilapisi lilin. Sebuah
\VUgras{tirai pintu} \textsuperscript{(22)} dari kain menutup cukup
rapat pintu yang menuju beranda. Di sanalah keluarga itu akan minum
kopi sesudah makan malam. \VUgras{Cangkir-cangkir}
\textsuperscript{(24)} dan \VUgras{lepek-lepek} \textsuperscript{(25)}
sudah siap di atas \VUgras{lemari piring} \textsuperscript{(23)}.

%%K t06-c3-06-1 p
6. --- Di atas meja, terpasang pada langit-langit, ada \VUgras{lampu
gantung} \textsuperscript{(26)} yang lampunya sangat terang dan pada
malam hari menerangi seluruh ruang makan.

%%K t06-c3-07-1 p
7. --- Sekarang mari kita jelajahi \VUgras{meja makan}
\textsuperscript{(27)} itu sendiri. Ketika keluarga itu masuk, alat
makan sudah siap. Di atas \VUgras{taplak meja} \textsuperscript{(28)}
telah diletakkan, pada tempat setiap orang yang makan, sebuah
\VUgras{piring} \textsuperscript{(33)}, sebuah \VUgras{serbet}
\textsuperscript{(29)}, sebuah \VUgras{sendok} \textsuperscript{(34)},
sebuah \VUgras{garpu} \textsuperscript{(35)}, sebuah \VUgras{pisau}
\textsuperscript{(36)} dan sebuah \VUgras{gelas} \textsuperscript{(32)}.
Ada pula \VUgras{botol-botol} \textsuperscript{(30)} untuk anggur dan
sebuah \VUgras{botol air} \textsuperscript{(31)}.

%%K t06-c3-08-1 p
8. --- \VUgras{Pelayan} \textsuperscript{(39)} mula-mula membawa
\VUgras{mangkuk sup} \textsuperscript{(37)}, yang sedikit supnya masih
berasap. Sekarang ia menyajikan \VUgras{hidangan}
\textsuperscript{(38)} pertama, hidangan daging. Kemudian ia akan
menyajikan yang sudah kita sebutkan tadi ; akhirnya, sesudah
\VUgras{hidangan penutup,} ia akan memanfaatkan saat tuan-tuannya
masuk ke beranda untuk mengangkat alat-alat meja dan merapikan kembali
ruang makan.

%%K t06-c3-08-2 p
\textit{Catatan.} --- \textit{Kata-kata sehari-hari yang berkaitan
dengan makanan hanya disebutkan secara ringkas di sini. Daftarnya akan
dilengkapi pada kedua bagan « Hotel » (no. 13) dan « Pasar » (no.
15).}

%%K t06-tit-7 sub
\VUsaut{4.6mm}
\VUpk{10.2pt}{40}{Dapur.}
\VUsaut{3.0mm}

%%K t06-c4-01-1 p
1. --- Di dapur, yang mulai kita lihat lewat pintu yang setengah
terbuka, kita melihat kekacauan yang menarik. Di depan \VUgras{tungku}
\textsuperscript{(1)} tempat \VUgras{api} \textsuperscript{(3)}
berderak, terpasang sebuah \VUgras{pemutar panggangan}
\textsuperscript{(5)}. \VUgras{Daging panggang} \textsuperscript{(7)}
yang bagus \VUgras{ditusuk} \textsuperscript{(6)} dan sedang masak.

%%K t06-c4-02-1 p
2. --- \VUgras{Peniup api} \textsuperscript{(8)} dan \VUgras{sekop
arang} \textsuperscript{(9)}, yang baru saja dipakai, tertinggal di
lantai bersama \VUgras{kayu bakar} \textsuperscript{(10)}.

%%K t06-c4-03-1 p
3. --- Bagian atas perapian sudah dirapikan : \VUgras{lentera}
\textsuperscript{(11)}, \VUgras{tempat korek api} \textsuperscript{(13)}
yang berisi \VUgras{korek api} \textsuperscript{(12)}, \VUgras{kaki
lilin} \textsuperscript{(14)} dan \VUgras{tempat lilin}
\textsuperscript{(15)} dengan \VUgras{lilinnya} \textsuperscript{(16)}
sudah pada tempatnya. Tetapi \VUgras{ember arang}
\textsuperscript{(18)} yang besar, penuh \VUgras{batu bara}
\textsuperscript{(17)} yang baru saja diisi \VUgras{juru masak}
\textsuperscript{(23)} di ruang bawah tanah, tertinggal di tengah
dapur.

%%K t06-c4-04-1 p
4. --- Sungguh, perempuan yang malang itu harus bergegas supaya makan
siang siap tepat pada waktunya. Sekarang ia berada di dekat
\VUgras{kompor} \textsuperscript{(19)} dan mengaduk \VUgras{saus}
\textsuperscript{(24)} yang tidak boleh terlalu gosong.

%%K t06-c4-05-1 p
5. --- Di dalam \VUgras{oven} \textsuperscript{(20)} sedang memanggang
kue yang ia siapkan untuk makanan kecil anak-anak. Di atas kompor
tergantung \VUgras{sendok kuah} \textsuperscript{(21)} dan
\VUgras{sendok saring} \textsuperscript{(22)}.

%%K t06-tit-8 sub
\VUsaut{4.0mm}
\VUpk{10.2pt}{40}{Perkakas.}
\VUsaut{3.0mm}

%%K t06-c4-06-1 p
6. --- \VUgras{Perkakas masak} \textsuperscript{(25)} yang tergantung
di ujung ruangan sangat bersih. \VUgras{Panci-panci}
\textsuperscript{(26)} tembaga yang indah, \VUgras{tempat susu}
\textsuperscript{(27)}, \VUgras{saringan kecil} \textsuperscript{(28)}
dan \VUgras{parutan} \textsuperscript{(29)} telah dibersihkan dengan
saksama.

%%K t06-c4-07-1 p
7. --- \VUgras{Tempat garam} \textsuperscript{(30)} diletakkan di atas
lemari piring di dekat \VUgras{teko} \textsuperscript{(31)} dan
\VUgras{botol minyak besar} \textsuperscript{(32)}. \VUgras{Laci}
\textsuperscript{(33)} itu barangkali berisi alat makan dan pisau
dapur.

%%K t06-c4-08-1 p
8. --- \VUgras{Sapu} \textsuperscript{(34)} dan \VUgras{kemoceng}
\textsuperscript{(35)} ada di sudut. Juru masak baru saja memakai
\VUgras{talenan kayu} \textsuperscript{(36)} ; ia meninggalkannya di
dekat jendela.

%%K t06-c4-09-1 p
9. --- Sehelai \VUgras{lap tangan} \textsuperscript{(37)} tergantung
di dinding, dekat \VUgras{corong} \textsuperscript{(38)}. Di bagian
dapur itu disimpan \VUgras{pemanggang} \textsuperscript{(39)},
\VUgras{pisau cincang} \textsuperscript{(40)} dan \VUgras{talenan}
\textsuperscript{(41)}. Kemudian, di atas pisau cincang, pada sebuah
\VUgras{papan} \textsuperscript{(42)}, ada \VUgras{periuk-periuk}
\textsuperscript{(43)}, \VUgras{panci rebus} \textsuperscript{(44)}
dengan \VUgras{tutupnya} \textsuperscript{(45)}, dan \VUgras{kuali}
\textsuperscript{(46)}. \VUgras{Wajan} \textsuperscript{(47)}
tergantung di dekatnya.

%%K t06-c4-10-1 p
10. --- Hanya \VUgras{keran} \textsuperscript{(48)} tembaga yang
memancarkan kilau di sudut ini. \VUgras{Bak cuci}
\textsuperscript{(49)}, yang di atasnya diletakkan \VUgras{tempayan}
\textsuperscript{(50)} yang tebal, tentu sangat praktis dengan lemari
kecilnya, tempat orang menyimpan \VUgras{tong cuka}
\textsuperscript{(51)} yang berkeran \textsuperscript{(52)},
\VUgras{tempat sampah} \textsuperscript{(53)} dan \VUgras{piring-piring
kotor} \textsuperscript{(54)}.

%%K t06-tit-9 sub
\VUsaut{4.0mm}
\VUpk{10.2pt}{40}{Barang-barang lain di dapur.}
\VUsaut{3.0mm}

%%K t06-c4-11-1 p
11. --- \VUgras{Bak besar} \textsuperscript{(55)} tempat orang mencuci
\VUgras{sayuran} \textsuperscript{(58)} dan \VUgras{kuali besi}
\textsuperscript{(56)} yang diletakkan di atas \VUgras{lantai ubin}
\textsuperscript{(57)} barangkali juga mempunyai tempatnya di dalam
lemari itu.

%%K t06-c4-12-1 p
12. --- Juru masak itu tentu tidak kekurangan tungku, sebab inilah
lagi, di sudut meja, sebuah \VUgras{kompor gas} \textsuperscript{(59)}
tempat air memanas di dalam panci kecil. \VUgras{Uap}
\textsuperscript{(60)} yang keluar darinya menunjukkan kepada kita
bahwa air itu barangkali mendidih. Boleh jadi air itu akan dipakai
untuk kopi, sebab inilah di ujung meja yang lain \VUgras{teko kopi}
\textsuperscript{(64)} dan \VUgras{penggiling kopi}
\textsuperscript{(63)}.

%%K t06-c4-13-1 p
13. --- Kompor itu dihubungkan dengan \VUgras{pembakar gas}
\textsuperscript{(62)} oleh sebuah \VUgras{selang karet}
\textsuperscript{(61)} yang panjang.

%%K t06-c4-14-1 p
14. --- \VUgras{Pembantu harian} \textsuperscript{(66)} yang datang
menolong juru masak sedang menyiapkan sayuran untuk makan malam. Bila
sayuran itu sudah dibersihkan dan dicuci, ia akan meletakkannya di
dalam \VUgras{baskom} \textsuperscript{(65)} sambil menunggu waktu
memasaknya.

%%K t06-tit-10 sub
\VUsaut{4.0mm}
{\centering\fontsize{10.2pt}{10.2pt}\selectfont\textls[40]{\scshape Kamar Mandi.} \textit{(Ruang mandi.)}\par}
\VUsaut{3.0mm}

%%K t06-c5-01-1 p
1. --- Tinggal satu kelancangan lagi yang harus kita lakukan untuk
mengenal kamar-kamar utama rumah itu : melihat perlengkapan kamar
mandi. Itu tidak akan memakan waktu lama, sebab kamar itu tidak besar,
dan kita akan segera tahu bahwa \VUgras{bak rendam}
\textsuperscript{(1)} mempunyai \VUgras{pemanas air}
\textsuperscript{(2)} yang baik, bahwa \VUgras{tempat sabun}
\textsuperscript{(3)} dapat dijangkau tangan orang yang mandi, dan
bahwa \VUgras{gantungan handuk} \textsuperscript{(5)} tentu memudahkan
mengeringkan \VUgras{handuk-handuk} \textsuperscript{(4)}.

%%K t06-c5-02-1 p
2. --- Dinding kamar itu ditutup \VUgras{ubin porselen}
\textsuperscript{(6)}, yang memungkinkan memasang \VUgras{pancuran
mandi} \textsuperscript{(7)} tanpa perlu khawatir bahwa air, yang
memercik kembali ketika dipakai, akan merusak tembok. Sebuah
\VUgras{baskom kecil} \textsuperscript{(8)} dari seng, yang berguna
untuk merendam kaki, melengkapi perabot itu.

\VUsaut{6.0mm}
\VUfilet{22mm}
